\documentclass[12pt,a4paper]{article}

% ------------------------------------------------
% Basis-Pakete (robust für Overleaf / pdfLaTeX)
% ------------------------------------------------
\usepackage[T1]{fontenc}
\usepackage[utf8]{inputenc}
\usepackage[ngerman]{babel}
\usepackage[a4paper,margin=2.5cm]{geometry}
\usepackage{setspace}
\usepackage{mathptmx}      % Times-Schrift, meist robust
\usepackage{graphicx}
\usepackage{hyperref}
\usepackage{longtable}
\usepackage{array}
\usepackage{booktabs}
\usepackage{enumitem}
\usepackage{fancyhdr}

% ------------------------------------------------
% Formatierung
% ------------------------------------------------
\onehalfspacing
\setlength{\parindent}{0pt}
\setlength{\parskip}{0.6em}

\hypersetup{
    colorlinks=true,
    linkcolor=black,
    urlcolor=blue,
    pdftitle={Praktikumsjournal},
    pdfauthor={Vorname Nachname}
}

\pagestyle{fancy}
\fancyhf{}
\fancyhead[L]{Praktikumsjournal}
\fancyhead[R]{\leftmark}
\fancyfoot[C]{\thepage}

% ------------------------------------------------
% Eigene Angaben
% ------------------------------------------------
\newcommand{\Name}{Loris Keller}
\newcommand{\Fach}{Physik}
\newcommand{\Adresse}{Rychenbergstrasse 140, 8400 Winterthur}
\newcommand{\Email}{Loris.keller@uzh.ch}
\newcommand{\Datum}{\today}
\newcommand{\Praktikum}{Unterrichtspraktikum I}
\newcommand{\Schule}{Kantonsschule Im Lee}
\newcommand{\Lehrperson}{Sara Romer}

\begin{document}

% =================================================
% TITELSEITE
% =================================================
\begin{titlepage}
    \thispagestyle{empty}
    
    % ---------- Logo oben rechts ----------
    \begin{flushright}
        % Falls du ein Logo hochlädst, z.B. uzh_logo.png, dann ersetze die nächste Zeile durch:
        \includegraphics[width=4cm]{UZH_logo.png}
    \end{flushright}

    \vspace{1.5cm}

    \begin{center}
        {\Large Universität Zürich\par}
        \vspace{0.3cm}
        {\large Institut für Erziehungswissenschaft\par}
        \vspace{0.3cm}
        {\large Lehrerinnen- und Lehrerbildung Maturitätsschulen\par}

        \vspace{2.5cm}

        {\Huge\bfseries Praktikumsjournal\par}
        \vspace{0.6cm}
        {\Large \Praktikum\par}

        \vspace{2cm}
    \end{center}

    \renewcommand{\arraystretch}{1.3}
    \begin{tabular}{>{\bfseries}p{4.5cm}p{9cm}}
        Name: & \Name \\
        Fach: & \Fach \\
        Schule: & \Schule \\
        Praktikumslehrperson: & \Lehrperson \\
        Adresse: & \Adresse \\
        E-Mail: & \Email \\
        Datum: & \Datum \\
    \end{tabular}

    \vfill
\end{titlepage}

% =================================================
% INHALTSVERZEICHNIS
% =================================================
\tableofcontents
\newpage

% =================================================
% EINLEITUNG
% =================================================
\section{Einleitung}

Dieses Praktikumsjournal dokumentiert und reflektiert zentrale Erfahrungen aus dem \Praktikum. 
Im Zentrum stehen zwei bis drei thematische Schwerpunkte beziehungsweise Fragestellungen, anhand 
derer eigene Unterrichtserfahrungen, Beobachtungen und Entwicklungsmöglichkeiten analysiert werden.

\subsection{Ausgangslage und Zielsetzung}

Beschreibe hier kurz:
\begin{itemize}[leftmargin=1.5cm]
    \item den Kontext des Praktikums,
    \item die Lerngruppe oder Klasse,
    \item deine Rolle im Praktikum,
    \item deine Ziele für das Journal.
\end{itemize}

\subsection{Gewählte thematische Schwerpunkte}

Formuliere hier zwei bis drei zentrale Fragestellungen.

Beispiel:
\begin{enumerate}[leftmargin=1.5cm]
    \item Wie gelang mir die Strukturierung von Unterricht?
    \item Inwiefern konnte ich Lernprozesse der Schülerinnen und Schüler unterstützen?
    \item Welche Herausforderungen zeigten sich im Bereich Klassenführung?
\end{enumerate}

% =================================================
% REFLEXION
% =================================================
\section{Reflexion}

\subsection{Thematischer Schwerpunkt 1: [Titel einfügen]}

\subsubsection*{Fragestellung}
Formuliere hier die konkrete Fragestellung.

\subsubsection*{Beschreibung der Situation}
Beschreibe eine konkrete Unterrichtssituation, Beobachtung, Planung oder ein Material, auf das du dich beziehst.

\subsubsection*{Analyse}
Analysiere die Situation:
\begin{itemize}[leftmargin=1.5cm]
    \item Was ist gelungen?
    \item Wo zeigten sich Schwierigkeiten?
    \item Welche Ursachen könnten dafür verantwortlich gewesen sein?
    \item Welche Schlüsse ziehst du daraus?
\end{itemize}

\subsubsection*{Bezug zu Theorie und Literatur}
Hier kannst du Bezüge zu allgemeindidaktischen, fachdidaktischen oder pädagogisch-psychologischen Konzepten herstellen.

\subsubsection*{Zwischenfazit}
Fasse die wichtigsten Erkenntnisse kurz zusammen.

\subsection{Thematischer Schwerpunkt 2: [Titel einfügen]}

\subsubsection*{Fragestellung}
Formuliere hier die konkrete Fragestellung.

\subsubsection*{Beschreibung der Situation}
Beschreibe die relevante Unterrichtssituation oder das Material.

\subsubsection*{Analyse}
Reflektiere die Situation vertieft und begründet.

\subsubsection*{Bezug zu Theorie und Literatur}
Verknüpfe deine Beobachtungen mit Literatur oder theoretischen Modellen.

\subsubsection*{Zwischenfazit}
Fasse die wichtigsten Erkenntnisse zusammen.

\subsection{Thematischer Schwerpunkt 3: [optional]}

Diesen Abschnitt kannst du verwenden, falls du drei Fragestellungen bearbeitest.

\subsubsection*{Fragestellung}
Formuliere hier die Fragestellung.

\subsubsection*{Beschreibung der Situation}
Beschreibe die relevante Praxissituation.

\subsubsection*{Analyse}
Reflektiere die Situation.

\subsubsection*{Bezug zu Theorie und Literatur}
Ordne deine Beobachtungen theoretisch ein.

\subsubsection*{Zwischenfazit}
Formuliere die wichtigsten Erkenntnisse.

% =================================================
% STANDORTBESTIMMUNG
% =================================================
\section{Standortbestimmung mit Entwicklungszielen}

\subsection{Bereits vorhandene Kompetenzen}

Reflektiere hier, welche Kompetenzen bereits sichtbar geworden sind.

Beispiele:
\begin{itemize}[leftmargin=1.5cm]
    \item fachliche Sicherheit,
    \item verständliche Erklärungen,
    \item klare Strukturierung,
    \item lernförderliche Interaktion,
    \item Klassenführung,
    \item reflektierter Umgang mit Rückmeldungen.
\end{itemize}

\subsection{Entwicklungsziele für das nächste Praktikum}

Formuliere zwei bis vier konkrete Entwicklungsziele.

\renewcommand{\arraystretch}{1.3}
\begin{longtable}{p{0.28\textwidth}p{0.30\textwidth}p{0.34\textwidth}}
\toprule
\textbf{Entwicklungsziel} & \textbf{Begründung} & \textbf{Konkrete Umsetzung} \\
\midrule
Ziel 1 & Warum ist dieses Ziel wichtig? & Welche konkreten Schritte planst du? \\
Ziel 2 & Warum ist dieses Ziel wichtig? & Welche konkreten Schritte planst du? \\
Ziel 3 & Optional & Optional \\
Ziel 4 & Optional & Optional \\
\bottomrule
\end{longtable}

\subsection{Persönliche Schlussreflexion}

Formuliere ein kurzes Gesamtfazit:
\begin{itemize}[leftmargin=1.5cm]
    \item Was nimmst du aus dem Praktikum mit?
    \item Welche Entwicklung ist bereits sichtbar?
    \item Worauf willst du im nächsten Praktikum besonders achten?
\end{itemize}

% =================================================
% ANHANG
% =================================================
\newpage
\section{Anhang}

Hier können ausgewählte Materialien aus dem Praktikum dokumentiert werden, zum Beispiel:
\begin{itemize}[leftmargin=1.5cm]
    \item Unterrichtsplanungen,
    \item Arbeitsblätter,
    \item Beobachtungsnotizen,
    \item Schülerarbeiten,
    \item Tafelbilder oder Präsentationsfolien.
\end{itemize}

\subsection{Übersicht der Anhangsdokumente}

\begin{enumerate}[leftmargin=1.5cm]
    \item Anhang A: Unterrichtsplanung vom [Datum]
    \item Anhang B: Arbeitsblatt zum Thema [Thema]
    \item Anhang C: Beobachtungsprotokoll
    \item Anhang D: Beispiel einer Schülerarbeit
\end{enumerate}

% =================================================
% LITERATURVERZEICHNIS
% =================================================
\newpage
\section*{Literaturverzeichnis}
\addcontentsline{toc}{section}{Literaturverzeichnis}

% Literatur hier von Hand eintragen, z.B.:
Autor, A. A. (Jahr). \textit{Titel}. Verlag.

\end{document}